% !TeX root = ../lista.tex
\begin{equation}
\frac{dN(t)}{dt} = rN(t) \tag{Eq 5.1}
\end{equation}
Cuja solução geral, dado um capital inicial $C_0$, é expressa como:
\begin{equation}
N(t) = C_0 e^{rt} \tag{Eq 5.4}
\end{equation}
Onde:
\begin{itemize}
  \item $N(t)$ é o valor acumulado no tempo $t$ (\$1.000.000).
  \item $C_0$ é o investimento inicial necessário no ano de 2002.
  \item $r$ é a taxa de juros anual contínua (5,5\% = 0,055).
  \item $t$ é o intervalo de tempo decorrido desde o investimento inicial.
\end{itemize}
Para determinar o valor a ser reservado ($C_0$), isolamos a variável na função inversa:
\begin{equation}
C_0 = N(t)\, e^{-rt}
\end{equation}

\paragraph{Caso 1 -- Projeção para o ano 2022 ($t=20$ anos):} considerando o intervalo $t=2022-2002=20$, aplicamos os parâmetros:

$C_0 = 1{.}000{.}000 \cdot e^{-(0{,}055 \cdot 20)}$

$C_0 = 1{.}000{.}000 \cdot e^{-1{,}1}$

$C_0 \approx 1{.}000{.}000 \cdot 0{,}332871$

$\mathbf{C_0 \approx \$332.871{,}08}$

\paragraph{Caso 2 -- Projeção para o ano 2042 ($t=40$ anos):} considerando o intervalo $t=2042-2002=40$:

$C_0 = 1{.}000{.}000 \cdot e^{-(0{,}055 \cdot 40)}$

$C_0 = 1{.}000{.}000 \cdot e^{-2{,}2}$

$C_0 \approx 1{.}000{.}000 \cdot 0{,}110803$

$\mathbf{C_0 \approx \$110.803{,}16}$